\documentclass[11pt,aspectratio=169]{beamer}
\usetheme{SimplePlus}

\usepackage[T1]{fontenc}
\usepackage[utf8]{inputenc}
\usepackage{amsmath,amssymb,mathtools}
\usepackage{booktabs}
\usepackage{array}
\usepackage{appendixnumberbeamer}
\usepackage{hyperref}

\newcommand{\E}{\mathbb{E}}
\newcommand{\Var}{\mathrm{Var}}
\newcommand{\Cov}{\mathrm{Cov}}
\newcommand{\Prob}{\mathbb{P}}
\newcommand{\ATT}{\mathrm{ATT}}
\newcommand{\RMSPE}{\mathrm{RMSPE}}
\newcommand{\tr}{\mathrm{tr}}
\newcommand{\indep}{\perp\!\!\!\perp}
\newcommand{\argmin}{\mathop{\mathrm{arg\,min}}}
\newcommand{\derivbutton}[1]{\vspace{0.4em}\hfill\hyperlink{app:#1}{\beamerbutton{Appendix details}}}
\newcommand{\backbutton}[1]{\vspace{0.4em}\hfill\hyperlink{main:#1}{\beamerbutton{Back to main slide}}}

\title{Applied Econometrics: Session 8}
\subtitle{Synthetic Control}
\author{PhD Intensive Course}
\institute{Lecture 8}
\date{\today}

\begin{document}

\begin{frame}
  \titlepage
\end{frame}

\begin{frame}{How These Slides Are Structured}
  \begin{itemize}
    \item The lecture is about one/few treated units observed over time.
    \item We build the missing counterfactual as a weighted average of untreated units.
    \item Each concept appears in words, a formula, and an applied interpretation.
    \item Technical derivations and inference details are in the appendix.
  \end{itemize}
\end{frame}

\begin{frame}{Roadmap for Today}
  \tableofcontents
\end{frame}

\section{Motivation}

\begin{frame}{Why Synthetic Control?}
  \begin{itemize}
    \item Many important policies affect one state, one city, one school, or one firm.
    \item We may observe many years, but only one treated unit.
    \item A single control unit is rarely credible enough.
    \item Synthetic control constructs a comparison unit from many untreated units.
  \end{itemize}
  \begin{block}{Core idea}
    Instead of searching for the perfect untreated unit, build one from a weighted average of imperfect untreated units.
  \end{block}
\end{frame}

\begin{frame}{Running Example}
  Imagine a city introduces a congestion charge in year \(T_0+1\).
  \begin{itemize}
    \item Treated unit: the city that adopts the policy.
    \item Outcome: traffic delays, pollution, retail sales, or public transit use.
    \item Donor pool: similar cities that did not adopt the policy.
    \item Goal: estimate what would have happened in the treated city without the policy.
  \end{itemize}
  \begin{block}{Why not just compare to one city?}
    The closest city may match population but not income, match income but not trends, or match trends but not baseline levels.
  \end{block}
\end{frame}

\begin{frame}{The DiD Problem with One Treated Unit}
  Difference-in-differences usually asks for a treated group and a control group.
  \vspace{0.3em}

  With city-level data:
  \[
    \widehat{\tau}_{DiD}
    =
    (Y_{1,post}-Y_{1,pre})-(Y_{0,post}-Y_{0,pre}).
  \]
  \begin{itemize}
    \item The formula is simple, but the control unit matters enormously.
    \item With one treated and one control city, we cannot average away idiosyncratic shocks.
    \item Standard regression standard errors can be meaningless when the effective number of units is tiny.
  \end{itemize}
\end{frame}

\begin{frame}{From One Control to a Synthetic Control}
  Instead of comparing city \(1\) to city \(2\), compare city \(1\) to
  \[
    \text{synthetic city}
    =
    0.35 \times \text{city A}
    +
    0.25 \times \text{city B}
    +
    0.40 \times \text{city C}.
  \]
  \begin{itemize}
    \item The weights are chosen using only pre-treatment information.
    \item A good synthetic city should track the treated city before treatment.
    \item After treatment, the gap is interpreted as the estimated policy effect.
  \end{itemize}
\end{frame}

\begin{frame}{What Makes the Method Attractive?}
  \begin{itemize}
    \item Transparent: students can see which donor units receive weight.
    \item Visual: the treated path and synthetic path can be plotted over time.
    \item Design-based: credibility comes from pre-treatment fit and donor-pool logic.
    \item Useful when treatment is aggregate and cannot be studied with many treated units.
  \end{itemize}
  \begin{block}{Teaching warning}
    Synthetic control is not magic. It replaces a hidden extrapolation with a visible construction that we can inspect.
  \end{block}
\end{frame}

\section{Setup}

\begin{frame}{Data Structure}
  \hypertarget{main:setup}{}
  We observe \(J+1\) units over \(T\) periods.
  \begin{itemize}
    \item Unit \(1\) is treated after period \(T_0\).
    \item Units \(2,\ldots,J+1\) are untreated donor units.
    \item \(Y_{jt}\) is the observed outcome for unit \(j\) at time \(t\).
  \end{itemize}
  Potential outcomes:
  \[
    Y^N_{jt} = \text{outcome without treatment}, \qquad
    Y^I_{jt} = \text{outcome with treatment}.
  \]
  For the treated unit and \(t>T_0\),
  \[
    \tau_{1t}=Y^I_{1t}-Y^N_{1t}.
  \]
  \derivbutton{potential-outcomes}
\end{frame}

\begin{frame}{The Missing Counterfactual}
  For \(t>T_0\), the treated unit's observed outcome is
  \[
    Y_{1t}=Y^I_{1t}.
  \]
  The unobserved object is
  \[
    Y^N_{1t},
  \]
  the outcome the same treated unit would have had without treatment.
  \begin{block}{Synthetic-control target}
    Estimate \(Y^N_{1t}\) using untreated units that resemble the treated unit before treatment.
  \end{block}
\end{frame}

\begin{frame}{Donor Pool}
  The donor pool is the set of units allowed to contribute to the synthetic control.
  \vspace{0.3em}

  A good donor unit should:
  \begin{itemize}
    \item not receive the treatment;
    \item not be strongly affected by spillovers from the treated unit;
    \item be measured consistently over time;
    \item be plausibly comparable in the pre-treatment period.
  \end{itemize}
  \begin{block}{Practical rule}
    The donor pool is a research-design choice, not a software default.
  \end{block}
\end{frame}

\begin{frame}{Bad Donors Contaminate the Counterfactual}
  Do not include donors that:
  \begin{itemize}
    \item adopt the same policy soon after the treated unit;
    \item respond strategically to the treated unit's policy;
    \item have data definitions that change at the intervention date;
    \item are structurally impossible comparisons.
  \end{itemize}
  \begin{block}{Example}
    If California adopts a tobacco tax, a state with a simultaneous anti-smoking campaign is not a clean donor for the no-policy path.
  \end{block}
\end{frame}

\begin{frame}{Pre-Treatment Information}
  Synthetic control chooses weights using predictors measured before treatment.
  Common predictors:
  \begin{itemize}
    \item lagged outcomes \(Y_{j1},\ldots,Y_{jT_0}\);
    \item baseline covariates such as income, demographics, or prices;
    \item averages over selected pre-treatment windows;
    \item other variables believed to forecast the outcome.
  \end{itemize}
  \begin{block}{Beginner warning}
    Post-treatment variables must not be used to choose weights. That would build the answer into the counterfactual.
  \end{block}
\end{frame}

\begin{frame}{Predictor Matrices}
  Let \(X_1\) be the \(k\times 1\) vector of pre-treatment predictors for the treated unit.
  \[
    X_1=
    \begin{bmatrix}
      \text{treated pre-period outcome average}\\
      \text{treated income}\\
      \text{treated baseline demographics}\\
      \vdots
    \end{bmatrix}.
  \]
  Let \(X_0\) be the \(k\times J\) matrix with the same predictors for donor units.
  \[
    X_0=
    \begin{bmatrix}
      \text{donor 1 predictor} & \cdots & \text{donor J predictor}
    \end{bmatrix}.
  \]
\end{frame}

\section{Constructing the Synthetic Unit}

\begin{frame}{Weighted Counterfactual}
  \hypertarget{main:weighted-counterfactual}{}
  Choose weights
  \[
    W=(w_2,\ldots,w_{J+1})'.
  \]
  The synthetic estimate of the untreated outcome path is
  \[
    \widehat{Y}^N_{1t}
    =
    \sum_{j=2}^{J+1} w_jY_{jt}.
  \]
  The estimated treatment effect at time \(t>T_0\) is
  \[
    \widehat{\tau}_{1t}
    =
    Y_{1t}-\widehat{Y}^N_{1t}.
  \]
  \derivbutton{weighted-counterfactual}
\end{frame}

\begin{frame}{Convex Weights}
  Synthetic control usually imposes:
  \[
    w_j\geq 0 \quad \text{for all } j,
    \qquad
    \sum_{j=2}^{J+1}w_j=1.
  \]
  \begin{itemize}
    \item Nonnegative weights prevent using a donor with a negative contribution.
    \item Sum-to-one weights make the synthetic unit an average, not an arbitrary scale change.
    \item Together, these constraints keep the synthetic unit inside the donor-pool range.
  \end{itemize}
  \begin{block}{Intuition}
    The method is safer when it interpolates among observed units rather than extrapolating beyond them.
  \end{block}
  \derivbutton{convex-hull}
\end{frame}

\begin{frame}{The Optimization Problem}
  \hypertarget{main:optimization}{}
  A standard weight choice solves
  \[
    \widehat{W}
    =
    \argmin_W
    (X_1-X_0W)'V(X_1-X_0W)
  \]
  subject to
  \[
    w_j\geq 0,\qquad \sum_j w_j=1.
  \]
  \begin{itemize}
    \item \(V\) is a diagonal matrix that says which predictors matter more.
    \item With lagged outcomes only and \(V=I\), this is constrained least squares on the pre-treatment path.
    \item The fitted synthetic path is then carried into the post-treatment period.
  \end{itemize}
  \derivbutton{optimization}
\end{frame}

\begin{frame}{Why Not Plain OLS Weights?}
  If we regress the treated unit's pre-treatment outcome on donor outcomes, OLS may choose:
  \[
    \widehat{Y}_{1t}^{OLS}
    =
    -0.8Y_{2t}+1.7Y_{3t}+0.4Y_{4t}-0.3Y_{5t}.
  \]
  \begin{itemize}
    \item Negative weights can create impossible synthetic units.
    \item Weights above one can stretch donors far outside observed support.
    \item With many donor units, OLS can overfit the pre-period and behave wildly afterward.
  \end{itemize}
  \begin{block}{Practical lesson}
    A perfect pre-treatment fit from unconstrained OLS can be a warning sign, not a victory.
  \end{block}
\end{frame}

\begin{frame}{Small Numerical Illustration}
  Suppose the treated unit has pre-treatment predictors:
  \[
    X_1=(10, 4)'.
  \]
  Donors have
  \[
  \begin{array}{c|cc}
    \toprule
    \text{donor} & \text{income} & \text{baseline outcome}\\
    \midrule
    A & 8 & 3\\
    B & 12 & 5\\
    C & 9 & 4\\
    \bottomrule
  \end{array}
  \]
  One synthetic control is
  \[
    0.50 A + 0.25 B + 0.25 C
    =
    (9.25,3.75)'.
  \]
  It is not exact, but it is interpretable.
\end{frame}

\begin{frame}{Pre-Treatment Fit}
  \hypertarget{main:pretreatment-fit}{}
  \small
  The central diagnostic is:
  \[
    Y_{1t} \approx \sum_{j=2}^{J+1}\widehat{w}_jY_{jt}
    \qquad \text{for } t\leq T_0.
  \]
  \begin{itemize}
    \item If the synthetic unit cannot track the treated unit before treatment, the post-treatment comparison is weak.
    \item Pre-treatment fit should be shown graphically.
    \item A numeric fit measure is useful, but it cannot replace the plot.
  \end{itemize}
  \begin{block}{Fit measure}
    \[
      \RMSPE_{pre}
      =
      \sqrt{\frac{1}{T_0}\sum_{t=1}^{T_0}(Y_{1t}-\widehat{Y}^N_{1t})^2}.
    \]
  \end{block}
  \derivbutton{rmspe}
\end{frame}

\begin{frame}{What Counts as Good Fit?}
  There is no universal cutoff.
  \begin{itemize}
    \item Good fit means the synthetic path captures the level and trend of the treated path.
    \item Small errors in irrelevant wiggles matter less than persistent pre-treatment gaps.
    \item If the treated unit is outside the donor support, good fit may be impossible.
  \end{itemize}
  \begin{block}{Credibility statement}
    We trust the post-treatment gap only to the extent that the pre-treatment construction looked like a credible untreated version of the treated unit.
  \end{block}
\end{frame}

\begin{frame}{Treated vs. Synthetic Path}
  The main plot has two lines:
  \[
    \text{treated path } Y_{1t},
    \qquad
    \text{synthetic path } \widehat{Y}^N_{1t}.
  \]
  \begin{itemize}
    \item Before \(T_0\): the lines should be close.
    \item After \(T_0\): the vertical distance is the estimated effect.
    \item The treatment date should be marked clearly.
  \end{itemize}
  \begin{block}{Reading the graph}
    If the treated line falls below the synthetic line after a tax, the estimate says the tax reduced the outcome relative to the no-tax counterfactual.
  \end{block}
\end{frame}

\begin{frame}{The Effect Path}
  A second useful plot shows the gap:
  \[
    \widehat{\tau}_{1t}
    =
    Y_{1t}-\widehat{Y}^N_{1t}.
  \]
  \begin{itemize}
    \item This makes dynamic effects visible.
    \item Effects may grow, fade, reverse, or appear with delay.
    \item Pre-treatment gaps should hover near zero if fit is good.
  \end{itemize}
  \begin{block}{Interpretation}
    Synthetic control estimates a time path of treatment effects, not only one average effect.
  \end{block}
\end{frame}

\section{Inference and Placebos}

\begin{frame}{Why Usual Standard Errors Are Awkward}
  With one treated unit, the relevant sample size is not the number of rows in the panel.
  \begin{itemize}
    \item Repeating the same unit over many years does not create many independent treated units.
    \item Regression standard errors can understate uncertainty.
    \item We need inference that respects the small number of aggregate units.
  \end{itemize}
  \begin{block}{Alternative}
    Use placebo exercises: pretend each donor unit was treated and compare its fake effect with the real treated effect.
  \end{block}
\end{frame}

\begin{frame}{Placebo Test Logic}
  \hypertarget{main:placebo-logic}{}
  For each donor unit \(q\):
  \begin{enumerate}
    \item Temporarily label \(q\) as treated.
    \item Build a synthetic control for \(q\) using the remaining donors.
    \item Compute the placebo gap \(Y_{qt}-\widehat{Y}^N_{qt}\).
    \item Compare placebo gaps with the treated unit's gap.
  \end{enumerate}
  \begin{block}{Intuition}
    If many untreated units show gaps as large as the treated unit, the treated gap is less impressive.
  \end{block}
  \derivbutton{placebo}
\end{frame}

\begin{frame}{Graphical Placebo Inference}
  A common figure overlays:
  \begin{itemize}
    \item the treated unit's gap path, emphasized;
    \item all donor placebo gap paths, shown lightly;
    \item a vertical line at the treatment date;
    \item a horizontal zero line.
  \end{itemize}
  \begin{block}{Reading the figure}
    The treated effect is more persuasive when it becomes unusual after treatment but was not unusual before treatment.
  \end{block}
\end{frame}

\begin{frame}{Placebo \(p\)-Values}
  For a negative effect at final period \(T\), define
  \[
    p
    =
    \frac{
      1+\sum_{q\in\mathcal{D}}
      1\{\widehat{\tau}_{qT}^{placebo}\leq \widehat{\tau}_{1T}\}
    }{
      1+|\mathcal{D}|
    }.
  \]
  \begin{itemize}
    \item The numerator counts placebo effects at least as negative as the treated effect.
    \item The added \(1\) includes the treated unit in the randomization distribution.
    \item With few donors, \(p\)-values are necessarily coarse.
  \end{itemize}
  \derivbutton{placebo}
\end{frame}

\begin{frame}{Poorly Fit Placebos}
  Some placebo units cannot be matched well before treatment.
  \begin{itemize}
    \item Their post-treatment gaps are noisy even when there is no treatment.
    \item Comparing the treated unit to badly fitted placebos can be misleading.
    \item Researchers often report placebo graphs both with and without high pre-period RMSPE units.
  \end{itemize}
  \begin{block}{Transparent rule}
    If you exclude placebo units with poor pre-treatment fit, state the rule before interpreting the post-treatment gaps.
  \end{block}
\end{frame}

\begin{frame}{RMSPE Ratio}
  Another diagnostic statistic is
  \[
    \text{ratio}
    =
    \frac{\RMSPE_{post}}{\RMSPE_{pre}}.
  \]
  \begin{itemize}
    \item Large ratios mean the gap becomes much larger after treatment.
    \item This adjusts for units that were already hard to fit before treatment.
    \item It is useful for ranking the treated unit against placebo units.
  \end{itemize}
\end{frame}

\section{Interpretation and Workflow}

\begin{frame}{What Exactly Is Identified?}
  Synthetic control estimates the treated unit's missing untreated path:
  \[
    Y^N_{1t}
    \approx
    \widehat{Y}^N_{1t}
    =
    \sum_j\widehat{w}_jY_{jt}.
  \]
  Therefore
  \[
    \widehat{\tau}_{1t}=Y_{1t}-\widehat{Y}^N_{1t}
  \]
  is an estimate of the treatment effect for the treated unit at time \(t\).
  \begin{block}{Scope}
    This is not automatically the average effect for all cities, states, firms, or countries.
  \end{block}
\end{frame}

\begin{frame}{Substantive Interpretation}
  After estimating the gap, ask:
  \begin{itemize}
    \item Is the sign economically meaningful?
    \item Is the magnitude large relative to the baseline outcome?
    \item Does the timing match the policy mechanism?
    \item Are there other events at the treatment date that could explain the gap?
  \end{itemize}
  \begin{block}{Example}
    A tobacco-tax effect should not only be statistically unusual; it should also appear after the tax and have a plausible size given prices and consumption behavior.
  \end{block}
\end{frame}

\begin{frame}{Practical Workflow}
  \hypertarget{main:workflow}{}
  \begin{enumerate}
    \item Define the treated unit, treatment date, outcome, and estimand.
    \item Build a defensible donor pool.
    \item Choose pre-treatment predictors and scale them if needed.
    \item Solve for convex weights.
    \item Plot treated versus synthetic paths.
    \item Plot the gap path.
    \item Run placebo tests and report fit diagnostics.
    \item Interpret effect size, timing, and limitations.
  \end{enumerate}
  \derivbutton{workflow-checklist}
\end{frame}

\begin{frame}{Choosing Predictors}
  A simple starting point is to match lagged outcomes:
  \[
    X_1=(Y_{11},\ldots,Y_{1T_0})',
    \qquad
    X_0=
    \begin{bmatrix}
      Y_{21} & \cdots & Y_{J+1,1}\\
      \vdots & & \vdots\\
      Y_{2T_0} & \cdots & Y_{J+1,T_0}
    \end{bmatrix}.
  \]
  You may add covariates if they predict the outcome and are pre-treatment.
  \begin{block}{Practical warning}
    If one predictor is measured in thousands and another in percentages, scaling choices can dominate the weights.
  \end{block}
\end{frame}

\begin{frame}{Sensitivity Analysis}
  Report how conclusions change when:
  \begin{itemize}
    \item excluding one influential donor at a time;
    \item changing the pre-treatment window;
    \item changing predictor sets;
    \item dropping donors with possible spillovers;
    \item using alternative outcome definitions.
  \end{itemize}
  \begin{block}{Goal}
    We do not need every specification to be identical, but the main conclusion should not depend on one fragile choice.
  \end{block}
\end{frame}

\begin{frame}{Common Limitations}
  \begin{itemize}
    \item Poor pre-treatment fit weakens the design.
    \item Spillovers violate the donor-pool logic.
    \item Other shocks at the treatment date can be confounded with treatment.
    \item Convex weights cannot match treated units outside donor support.
    \item Many researcher choices remain: donors, predictors, windows, and fit thresholds.
  \end{itemize}
  \begin{block}{Bottom line}
    Synthetic control is most convincing when the design story, pre-treatment fit, and placebo evidence all point in the same direction.
  \end{block}
\end{frame}

\begin{frame}{Synthetic Control Compared with Earlier Methods}
  \begin{tabular}{p{0.24\linewidth}p{0.32\linewidth}p{0.32\linewidth}}
    \toprule
    \textbf{Method} & \textbf{Comparison} & \textbf{Key assumption}\\
    \midrule
    DiD & Treated group versus control group before/after & Parallel trends\\
    Matching & Treated units versus similar controls & Selection on observables\\
    Synthetic control & Treated unit versus weighted donor pool & Donor weights recover missing untreated path\\
    \bottomrule
  \end{tabular}
  \vspace{0.6em}

  \begin{block}{Connection}
    Synthetic control borrows matching's comparison logic and DiD's time-series logic, but makes the comparison unit explicit.
  \end{block}
\end{frame}

\begin{frame}{Python Example in the Course Folder}
  The companion script is:
  \[
    \texttt{lectures/code/lecture-8-synthetic-control.py}.
  \]
  It:
  \begin{itemize}
    \item simulates one treated unit and several donor units;
    \item estimates convex synthetic-control weights with projected gradient descent;
    \item compares treated and synthetic outcome paths;
    \item runs placebo reassignments across donor units;
    \item reports a finite-sample placebo \(p\)-value.
  \end{itemize}
  Run with plots:
  \[
    \texttt{python lectures/code/lecture-8-synthetic-control.py -{}-plot}
  \]
\end{frame}

\begin{frame}{What Students Should Be Able to Do}
  After this lecture, students should be able to:
  \begin{itemize}
    \item explain synthetic control as a missing-counterfactual estimator;
    \item define treated unit, donor pool, pre-treatment period, and effect path;
    \item write the convex-weight optimization problem;
    \item interpret treated-versus-synthetic and gap plots;
    \item explain placebo tests and pre-treatment fit diagnostics;
    \item name the main limitations and credibility checks.
  \end{itemize}
\end{frame}

\begin{frame}{Main Takeaways}
  \begin{itemize}
    \item Synthetic control builds a counterfactual unit from untreated donor units.
    \item Convex weights make the construction transparent and reduce extrapolation risk.
    \item Pre-treatment fit is the central diagnostic.
    \item The effect is a post-treatment gap between the treated and synthetic paths.
    \item Placebo tests ask whether the treated gap is unusual relative to fake treated units.
  \end{itemize}
\end{frame}

\begin{frame}{Sources and References}
  \begin{itemize}
    \item Matheus Facure Alves, \emph{Causal Inference for the Brave and True}, Chapter 15: Synthetic Control. \url{https://matheusfacure.github.io/python-causality-handbook/15-Synthetic-Control.html}
    \item Abadie, Diamond, and Hainmueller (2010), ``Synthetic Control Methods for Comparative Case Studies.'' \emph{Journal of the American Statistical Association}.
    \item Abadie (2021), ``Using Synthetic Controls: Feasibility, Data Requirements, and Methodological Aspects.'' \emph{Journal of Economic Literature}.
  \end{itemize}
\end{frame}

\appendix

\begin{frame}{Appendix Map}
  \begin{itemize}
    \item A1: Potential-outcomes setup for synthetic control.
    \item A2: Why weighted donor outcomes estimate the missing path.
    \item A3: Convex weights and the convex hull.
    \item A4: Constrained least-squares optimization.
    \item A5: RMSPE diagnostics.
    \item A6: Placebo inference logic.
    \item A7: Practical reporting checklist.
  \end{itemize}
\end{frame}

\begin{frame}{Appendix A1: Potential-Outcomes Setup}
  \hypertarget{app:potential-outcomes}{}
  For the treated unit \(1\):
  \[
    Y_{1t}=
    \begin{cases}
      Y^N_{1t}, & t\leq T_0,\\
      Y^I_{1t}, & t>T_0.
    \end{cases}
  \]
  For donor units \(j=2,\ldots,J+1\):
  \[
    Y_{jt}=Y^N_{jt}\quad \text{for all }t.
  \]
  The post-treatment effect is
  \[
    \tau_{1t}=Y^I_{1t}-Y^N_{1t},\qquad t>T_0.
  \]
  Since \(Y^I_{1t}\) is observed after treatment, identification reduces to approximating \(Y^N_{1t}\).
  \backbutton{setup}
\end{frame}

\begin{frame}{Appendix A2: Weighted Donor Counterfactual}
  \hypertarget{app:weighted-counterfactual}{}
  \small
  Suppose there exists a donor-weight vector \(W^\star\) such that the untreated treated-unit path is well approximated by
  \[
    Y^N_{1t}\approx \sum_{j=2}^{J+1}w^\star_jY^N_{jt}.
  \]
  Since donor units are untreated,
  \[
    Y^N_{jt}=Y_{jt}.
  \]
  Therefore
  \[
    Y^N_{1t}\approx \sum_{j=2}^{J+1}w^\star_jY_{jt}.
  \]
  Replacing \(W^\star\) with the estimated pre-treatment weights gives
  \[
    \widehat{\tau}_{1t}=Y_{1t}-\sum_{j=2}^{J+1}\widehat{w}_jY_{jt}.
  \]
  \backbutton{weighted-counterfactual}
\end{frame}

\begin{frame}{Appendix A3: Convex Hull Logic}
  \hypertarget{app:convex-hull}{}
  A convex combination has
  \[
    w_j\geq 0,\qquad \sum_jw_j=1.
  \]
  Let \(Z_j\) be any scalar predictor. The synthetic predictor is
  \[
    Z_{syn}=\sum_jw_jZ_j.
  \]
  If \(Z_{\min}\leq Z_j\leq Z_{\max}\) for all donors, then
  \[
    Z_{\min}
    =
    \sum_jw_jZ_{\min}
    \leq
    \sum_jw_jZ_j
    \leq
    \sum_jw_jZ_{\max}
    =
    Z_{\max}.
  \]
  So the synthetic unit cannot lie outside the donor range on any one predictor.
  \backbutton{optimization}
\end{frame}

\begin{frame}{Appendix A4: Constrained Least Squares}
  \hypertarget{app:optimization}{}
  With \(V=I\), the weight problem is
  \[
    \min_W \; (X_1-X_0W)'(X_1-X_0W)
  \]
  subject to
  \[
    \iota'W=1,\qquad W\geq 0.
  \]
  Ignoring inequality constraints for a moment, the Lagrangian is
  \[
    \mathcal{L}(W,\lambda)
    =
    (X_1-X_0W)'(X_1-X_0W)+\lambda(\iota'W-1).
  \]
  First-order conditions:
  \[
    -2X_0'(X_1-X_0W)+\lambda\iota=0,
    \qquad
    \iota'W=1.
  \]
  Nonnegativity adds complementary-slackness conditions and is usually handled numerically.
  \backbutton{optimization}
\end{frame}

\begin{frame}{Appendix A5: RMSPE Diagnostics}
  \hypertarget{app:rmspe}{}
  Pre-treatment RMSPE:
  \[
    \RMSPE_{pre}
    =
    \sqrt{
    \frac{1}{T_0}
    \sum_{t=1}^{T_0}
    (Y_{1t}-\widehat{Y}^{N}_{1t})^2
    }.
  \]
  Post-treatment RMSPE over periods \(T_0+1,\ldots,T\):
  \[
    \RMSPE_{post}
    =
    \sqrt{
    \frac{1}{T-T_0}
    \sum_{t=T_0+1}^{T}
    (Y_{1t}-\widehat{Y}^{N}_{1t})^2
    }.
  \]
  The ratio \(\RMSPE_{post}/\RMSPE_{pre}\) is large when the post-treatment gap is large relative to pre-treatment mismatch.
  \backbutton{pretreatment-fit}
\end{frame}

\begin{frame}{Appendix A6: Placebo Inference as Randomization Logic}
  \hypertarget{app:placebo}{}
  Let \(\mathcal{D}\) be donor units. For each \(q\in\mathcal{D}\), compute a placebo statistic \(S_q\), such as final-period gap or RMSPE ratio.
  \vspace{0.3em}

  If the treatment assignment were exchangeable across units, then the treated statistic \(S_1\) should look like one draw from
  \[
    \{S_1,S_q:q\in\mathcal{D}\}.
  \]
  For a lower-tail effect:
  \[
    p=
    \frac{
      1+\sum_{q\in\mathcal{D}}1\{S_q\leq S_1\}
    }{
      1+|\mathcal{D}|
    }.
  \]
  This is exact only under the maintained exchangeability/design assumptions.
  \backbutton{placebo-logic}
\end{frame}

\begin{frame}{Appendix A7: Reporting Checklist}
  \hypertarget{app:workflow-checklist}{}
  A complete synthetic-control application should report:
  \begin{itemize}
    \item treated unit, treatment date, outcome, and analysis window;
    \item donor-pool inclusion and exclusion rules;
    \item predictor variables and scaling choices;
    \item estimated donor weights;
    \item treated-versus-synthetic path plot;
    \item gap plot and effect magnitudes;
    \item pre-treatment RMSPE and placebo evidence;
    \item limitations: spillovers, other shocks, poor fit, and sensitivity.
  \end{itemize}
  \backbutton{workflow}
\end{frame}

\end{document}
